\documentclass[a4paper,11pt]{article}
\usepackage[margin=2.5cm]{geometry}
\usepackage{amsmath,amssymb,amsthm}
\usepackage{booktabs}
\usepackage{hyperref}
\usepackage{xcolor}
\usepackage{listings}
\usepackage{enumitem}

\hypersetup{colorlinks=true, linkcolor=blue!50!black, urlcolor=blue!60!black}

\lstset{
  basicstyle=\ttfamily\small,
  breaklines=true,
  frame=single,
  framesep=4pt,
  rulecolor=\color{gray!50},
  backgroundcolor=\color{gray!5},
  columns=fullflexible,
}

\newtheorem{theorem}{Theorem}
\newtheorem{lemma}{Lemma}
\newtheorem{remark}{Remark}

\title{Phase 3 Vandermonde Encoder: \\
       \large DFT-node power sums and sort-encodings of the
              rep alignment table}
\author{Echinocoder / symcoder working notes}
\date{2026-05-26}

\begin{document}
\maketitle

\begin{abstract}
This document describes the mathematics behind
\texttt{symcoder.encoders.phase3\_vandermonde}, an alternative Phase 3
one-shot whole-table encoder for the symcoder pipeline.  Like the
simplicial-complex Phase 3 encoder it consumes the alignment table $T$
of shape $(|G|, |\text{rep}|)$ and produces a permutation-invariant
embedding; unlike the simplicial encoder it does so via polynomial
projections (Vandermonde-power-sums) or via sorted-projections, both with
a one-line correctness proof.  The implementation uses complex
roots-of-unity nodes (a DFT setup) for the C-infinity variant to achieve
perfect numerical conditioning, and admits an optional dimensional
``scale'' parameter to fold out per-row unit mismatches.
\end{abstract}

\tableofcontents
\newpage

%% ------------------------------------------------------------------
\section{Setting and notation}

We embed a multiset of $n$ vectors in $\mathbb{R}^k$.  Concretely the
multiset is the column-set of the alignment table
\[
  T \in \mathbb{R}^{n \times k}, \qquad
  T_{i, j} \;=\; \text{eval}\!\bigl(g_i \cdot \text{rep\_atoms}[j],\ E\bigr),
\]
where $g_1, \ldots, g_n$ enumerate the group $G$ (so $n = |G|$) and
$\text{rep\_atoms}$ is the unpruned list of $|\text{rep}| = k$
canonical atoms over which Phase 3 encoders work.  We treat $T$ as a
multiset of $n$ rows of $\mathbb{R}^k$; the encoder must be invariant
under permutation of rows.

The label-permutation group $H$ acts on $T$ by left-translation of
$g$-index.  An $H$-invariant encoding therefore inherits multiset
invariance automatically.

%% ------------------------------------------------------------------
\section{Vandermonde projection}

For each ``Vandermonde node'' $t_c \in \mathbb{C}$ (one per
$c = 0, 1, \ldots, m_{\text{proj}} - 1$), define the linear functional
\[
  L_c \;:\; \mathbb{R}^k \to \mathbb{C}, \qquad
  L_c(x) \;=\; \sum_{j=1}^{k} t_c^{\,j-1}\, x_j.
\]
Apply $L_c$ to every row of $T$:
\[
  y_{c, i} \;=\; L_c(T_{i, \cdot}), \qquad i = 1, \ldots, n.
\]
Then $y_{c, \cdot}$ is a multiset (or rather a sequence
later treated as a multiset) of $n$ values in $\mathbb{C}$, one per row
of $T$.

%% ------------------------------------------------------------------
\section{Two ways to encode a multiset of $n$ scalars}

The encoder stores, for each $c$, an $n$-real representation of the
multiset $\{y_{c, 1}, \ldots, y_{c, n}\}$.  Two complementary choices
are implemented.

\subsection{Power-sum encoding (C-infinity mode)}

The classical Newton-Girard identities relate the power sums
$p_q = \sum_i y_i^q$ to the elementary symmetric polynomials $e_q$ of
$\{y_i\}$ via a triangular set of recursions.  Knowing $p_1, p_2, \ldots,
p_n$ therefore determines $e_1, \ldots, e_n$, which are the
coefficients (up to sign) of the polynomial
\[
  \prod_{i=1}^{n} (X - y_i)
  \;=\; X^n - e_1 X^{n-1} + e_2 X^{n-2} - \cdots \pm e_n,
\]
whose roots are precisely $\{y_i\}$.  Hence:

\begin{lemma}[Univariate Newton-Girard]
The first $n$ power sums $p_1, p_2, \ldots, p_n$ uniquely determine the
multiset of $n$ scalars $\{y_1, \ldots, y_n\}$.
\end{lemma}

The C-infinity encoding for projection $c$ is therefore
\[
  M_{c, p} \;=\; \sum_{i=1}^{n} y_{c, i}^{\,p}, \qquad p = 1, 2, \ldots, n.
\]
Each $M_{c, p}$ is a polynomial in the entries of $T$, so the encoding
is $C^\infty$ everywhere.

\subsection{Sort encoding (C-zero mode)}

The same multiset can equivalently be encoded by its ascending sort:
\[
  S_{c, p} \;=\; \text{sort}(y_{c, 1}, \ldots, y_{c, n})_p,
  \qquad p = 1, \ldots, n.
\]
Sorting is the canonical bijection between multisets of reals and
ascending tuples of reals.  $S_{c, p}$ is continuous and Lipschitz in
the entries of $T$, but not differentiable at value-ties.  It is
``dimensionally homogeneous'' with $T$ (no powers are taken), which is
why this mode is offered as an alternative to power sums.

\subsection{Bridging projections: how many nodes are enough?}

A multiset of $n$ values in $\mathbb{C}$ determined per projection
does not by itself recover the original multiset of $n$ vectors in
$\mathbb{R}^k$.  We need enough projections to disambiguate the joint
distribution.  Expanding the $p$-th power sum:
\[
  M_{c, p} \;=\; \sum_i \Bigl(\sum_j t_c^{j-1}\, T_{i, j}\Bigr)^{\!p}
  \;=\; \sum_{(a_1, \ldots, a_k):\,|a|=p}
        \binom{p}{a}\, t_c^{(j-1)\,a_j}\,
        \underbrace{\sum_i \prod_j T_{i, j}^{a_j}}_{\text{multi-symmetric
        power sum}\ P_{a}(T)}.
\]
So $M_{c, p}$ is a polynomial in $t_c$ of degree $(k-1)\,p$ whose
coefficients are the multivariate (multi-symmetric) power sums
$P_a(T)$ of total degree $|a| = p$.

\begin{theorem}
If $t_0, \ldots, t_{m-1}$ are $m$ distinct complex numbers with
$m \ge (k - 1)\,n + 1$, the values $\{M_{c, p} : 0 \le c < m,\
1 \le p \le n\}$ determine all multi-symmetric power sums of total
degree $\le n$ of the row-multiset of $T$.  Those in turn determine
the row-multiset itself.
\label{thm:bridging}
\end{theorem}

\begin{proof}[Proof sketch]
For each $p$, $M_{\cdot, p}$ as a function of $t$ is a polynomial of
degree $(k-1)\, p \le (k-1)\,n$ in $t$.  Knowing it at $(k-1)\,n + 1$
distinct points determines it (Vandermonde non-singularity).  Hence
all $P_a(T)$ for $|a| \le n$ are recovered.  The multi-symmetric
analogue of Newton-Girard (\emph{multi-symmetric power sums up to
total degree $n$ separate multisets of $n$ tuples}) then yields the
multiset of rows.
\end{proof}

The sort-encoding admits the same statement: per projection,
$\{S_{c, 1}, \ldots, S_{c, n}\}$ is equivalent in information content
to $\{M_{c, 1}, \ldots, M_{c, n}\}$ by Newton-Girard, and the
Vandermonde / multi-symmetric argument is identical.

%% ------------------------------------------------------------------
\section{Numerical issues and the choice of nodes}

Theorem~\ref{thm:bridging} requires only that the $t_c$ be distinct;
their magnitudes are otherwise free.  Two issues constrain the actual
choice in float64 arithmetic.

\subsection{Overflow with large-magnitude nodes}

The naive choice $t_c = c + 1$ (distinct positive integers) gives
$t_{m-1}^{\,k-1} \approx m^{k-1}$.  For the parity-fixture plan
($m = 289$, $k = 25$) this is $289^{24} \approx 10^{58}$, which
overflows float64 at the projection matrix stage before any sum is
even taken.

\subsection{Underflow with bounded-magnitude real nodes}

Putting $|t_c| < 1$ (e.g.\ Chebyshev nodes in $(-1, 1)$) fixes
overflow.  But $t_c^{\,k-1}$ then becomes very small for nodes near
zero, so those projection-matrix entries silently underflow to zero.
The Vandermonde matrix becomes ill-conditioned even when no element
literally overflows or underflows.

\subsection{Unit-circle (DFT) nodes resolve both}

The choice
\[
  t_c \;=\; \exp\!\Bigl(\frac{2\pi i\,c}{m_{\text{proj}}}\Bigr),
  \qquad c = 0, 1, \ldots, m_{\text{proj}} - 1
\]
gives the $m_{\text{proj}}$-th roots of unity.  All are distinct (so
Theorem~\ref{thm:bridging} applies).  All have modulus exactly $1$,
so $|t_c^{\,j-1}| = 1$ for every $j$: no overflow, no underflow.
Furthermore the resulting Vandermonde matrix is (up to a unitary
normalisation) the DFT matrix, which has condition number $1$
exactly.  This is the gold-standard numerical choice.

%% ------------------------------------------------------------------
\section{From complex outputs to a real-valued vector}

With complex $t_c$ the projections $y_{c, i}$ and power sums $M_{c, p}$
are complex.  To package the encoding as a real vector we exploit
conjugate symmetry.

\begin{lemma}
$M_p(t) = \sum_i (L_t(T_{i, \cdot}))^p$, viewed as a function of
$t \in \mathbb{C}$, is a polynomial in $t$ with real coefficients.
Therefore $M_p(\overline{t}) = \overline{M_p(t)}$.
\end{lemma}

\begin{proof}
$L_t(T_{i, \cdot}) = \sum_j t^{j-1} T_{i, j}$ with real $T_{i, j}$, so
the coefficient of $t^d$ in $L_t(T_{i, \cdot})$ is real.  Raising to
the $p$-th power and summing over $i$ preserves reality of
coefficients.
\end{proof}

\paragraph{Consequence.}
The lower-half roots of unity carry no additional information beyond
the upper-half ones: $M_p$ at $\overline{t}$ is determined by $M_p$ at
$t$.  So we store only the upper-half (and real-axis) nodes.  For
$m_{\text{proj}}$ odd there is exactly one real-axis node $t = +1$ at
$c = 0$; for $m_{\text{proj}}$ even there are two, $t = +1$ at
$c = 0$ and $t = -1$ at $c = m_{\text{proj}}/2$.  The encoding layout:

\begin{itemize}[nosep]
  \item $c = 0$ (node $t = +1$): emit $\mathrm{Re}(M_{0, p})$ for
        $p = 1, \ldots, n$.  ($n$ reals; $\mathrm{Im}$ is zero by the
        lemma.)
  \item For each $c = 1, \ldots, \lfloor (m_{\text{proj}} - 1)/2 \rfloor$
        (upper-half nodes): emit $\mathrm{Re}(M_{c, p})$ and
        $\mathrm{Im}(M_{c, p})$ interleaved, for $p = 1, \ldots, n$.
        ($2n$ reals per such $c$.)
  \item If $m_{\text{proj}}$ even, $c = m_{\text{proj}}/2$ (node
        $t = -1$): emit $\mathrm{Re}(M_{c, p})$ for $p = 1, \ldots, n$.
        ($n$ reals.)
\end{itemize}

Total reals:
\[
  n + 2n \cdot \left\lfloor \frac{m_{\text{proj}} - 1}{2}\right\rfloor
    + [m_{\text{proj}} \text{ even}] \cdot n
  \;=\; n \cdot m_{\text{proj}}.
\]
The same total size as if we had used $m_{\text{proj}}$ real-valued
nodes, but with perfect numerical conditioning.

%% ------------------------------------------------------------------
\section{The optional scale parameter}

\subsection{The unit-mismatch problem}

In the C-infinity mode the encoding $M_{c, p}$ has the same physical
units as $y^p$.  But the columns of $T$ themselves have different
units when the rep atoms are of different mass dimension: $\text{mag}$
columns are in units of length, $\text{dot}$ columns in length-squared,
$\varepsilon_3$ columns in length-cubed.  Mixing these in the
projection $L_c(T_{i, \cdot})$ and then raising to the $p$-th power
gives output magnitudes that span $\sim |y|^1$ to $\sim |y|^n$ across
$p$.  For $n = 12$ that is twelve orders of magnitude of dynamic range
baked into the encoding before anything else has happened.

\subsection{Dimensional homogenisation}

We define the \emph{mass dimension} of an operation as the homogeneity
degree of its $\text{eval\_fn}$ under uniform scaling of its vector
arguments:
\[
  \text{eval\_fn}(k v_1, \ldots, k v_r)
  \;=\; k^{\text{mass\_dim}}\, \text{eval\_fn}(v_1, \ldots, v_r).
\]
Concretely
$\text{mag}\!:1$,
$\text{magSq}\!:2$,
$\text{dot}\!:2$,
$\varepsilon_3\!:3$.
This degree is declared on the \texttt{Operation} object as the new
\texttt{mass\_dimension} field.

Given a user-supplied length scale $\lambda > 0$, the encoder first
divides each column of $T$ by $\lambda$ raised to that column's mass
dimension:
\[
  T'_{i, j} \;=\; T_{i, j}\, /\, \lambda^{\text{mass\_dim}(j)}.
\]
The resulting $T'$ is dimensionless and has entries of order unity
when $\lambda$ matches the natural scale of the event.  Power sums of
$T'$ stay within a few orders of magnitude regardless of $p$.

\paragraph{Why $\lambda$ must be external, not derived.}
A natural temptation is to set $\lambda$ to (say) the largest absolute
entry of $T$ at encoding time.  This is unsafe: the event values can
be arbitrarily close to zero (e.g.\ the parity-fixture event with
$p = (0,0,1), q = (0,0,-1)$ has many zero entries in $T$), so no
data-derived $\lambda$ is bounded away from zero.  A bad $\lambda$
amplifies rather than reduces dynamic range.

The scale is therefore a fixed encoder parameter and is recorded in
the segment descriptor so that downstream consumers can undo it.

%% ------------------------------------------------------------------
\section{Summary of guarantees and limitations}

\begin{itemize}[nosep]
  \item \textbf{Faithfulness.}  The encoder is deterministically
        faithful (separates distinct row-multisets of $T$) by
        Theorem~\ref{thm:bridging}.  No genericity assumption; no
        ``measure-zero failures''.
  \item \textbf{Smoothness.}  Cinf mode is $C^\infty$ in $T$
        (polynomials).  C0 mode is Lipschitz continuous but not
        differentiable at value-ties.
  \item \textbf{Numerical conditioning.}  Cinf mode uses DFT nodes,
        giving an implicit Vandermonde matrix of condition number $1$
        and no over/underflow in the projection step.
  \item \textbf{Output size.}  $m_{\text{proj}} \cdot n$ reals where
        $m_{\text{proj}} = (k-1)\,n + 1$, so $\mathcal{O}(n^2 k)$.
        Larger than the simplicial encoder's
        $2 n k + 1 - k = \mathcal{O}(nk)$ but with a one-line
        correctness proof and a smooth output (Cinf) or a
        dimensionally-homogeneous one (C0).
  \item \textbf{Scale parameter.}  In Cinf mode an external length
        scale $\lambda$ may be supplied to keep $M_{c, p}$ magnitudes
        bounded across $p$.  Doesn't affect faithfulness; only
        numerical conditioning.
  \item \textbf{What it does not do.}  Does not address the unit
        mismatch \emph{within} a single power $p$ across different
        $c$.  Does not provide an order-$O(nk)$ alternative to the
        simplicial encoder; that remains an open problem.
\end{itemize}

%% ------------------------------------------------------------------
\section{Code map}

\begin{tabular}{lp{9cm}}
\toprule
File & Role \\
\midrule
\texttt{symatom/atoms.py} &
  Adds the \texttt{mass\_dimension} field to \texttt{Operation}.
\\
\texttt{symcoder/operations/euclidean\{2,3\}.py} &
  Declares mass dimensions on the standard library operations.
\\
\texttt{symcoder/encoders/phase3\_common.py} &
  Shared helpers: \texttt{build\_alignment\_table},
  \texttt{column\_mass\_dimensions}, \texttt{apply\_scale}.
\\
\texttt{symcoder/encoders/phase3\_vandermonde.py} &
  This encoder.  \texttt{Phase3VandermondeEncoder} +
  \texttt{Phase3VandermondeEncoderFactory}.
\\
\texttt{symcoder/encoders/phase3\_factory.py} &
  Stacking wrapper.  \texttt{Phase3EncoderFactory} composes any
  number of Phase 3 sub-encoders.
\\
\texttt{symcoder/encoders/phase3\_simplicial.py} &
  The simplicial alternative.  Same input ($T$), different output style.
\\
\bottomrule
\end{tabular}

\end{document}
